\documentclass[11pt,a4paper]{article}

% ---------------------------------------------------------
% PREAMBLE (stable, warning-free, Overleaf-safe)
% ---------------------------------------------------------
\usepackage[utf8]{inputenc}
\usepackage[T1]{fontenc}
\usepackage{lmodern}
\usepackage{amsmath, amssymb, amsfonts}
\usepackage{bm}
\usepackage{geometry}
\usepackage{graphicx}
\usepackage{hyperref}
\usepackage{cite}
\usepackage{authblk}
\usepackage{physics}
\usepackage{mathtools}

\geometry{margin=2.4cm}

\title{\textbf{Companion LVIII: The Weak-Lensing Bispectrum in the Relaxation-Driven Cyclic Cosmology}}
\author{Michael Lehmann \\ \texttt{mi.lehmann@gmx.de}}
\date{April 2026}

\begin{document}
\maketitle

\begin{abstract}
The weak-lensing bispectrum probes the non-Gaussian structure of the gravitational 
potential by measuring the three-point correlations of the cosmic shear field. In 
the standard cosmological model, the bispectrum arises from nonlinear structure 
formation and the mode coupling induced by gravitational evolution. In the 
Relaxation-Driven Cyclic Cosmology (RDCC), the scale-dependent evolution of the 
gravitational potential modifies the nonlinear coupling, the potential depth, and 
the three-point structure of the shear field in a characteristic way. This 
Companion develops a continuous description of the weak-lensing bispectrum in 
RDCC, deriving how the relaxation scale influences the bispectrum amplitude, the 
configuration dependence, and the observational signatures in cosmic shear 
surveys. The resulting framework provides a distinctive probe of the RDCC 
gravitational field beyond the two-point level.
\end{abstract}
% ---------------------------------------------------------
\section{Introduction}

Weak gravitational lensing probes the integrated gravitational potential along the 
line of sight. While the power spectrum captures the Gaussian component of the 
shear field, the bispectrum captures the non-Gaussian structure induced by 
nonlinear gravitational evolution. The bispectrum is sensitive to the mode 
coupling between large and small scales and provides a powerful probe of the 
gravitational field.

In the standard cosmological model, the bispectrum arises from nonlinear structure 
formation and the coupling of density modes. In RDCC, the gravitational potential 
evolves differently. The relaxation mechanism suppresses the decay of small-scale 
modes and enhances the coherence of large-scale modes. This modifies the nonlinear 
coupling, the potential depth, and the three-point structure of the shear field in 
a scale-dependent manner, producing a distinctive signature in the weak-lensing 
bispectrum.
% ---------------------------------------------------------
\section{The RDCC Nonlinear Potential}

The nonlinear gravitational potential evolves as
\begin{equation}
    \Phi_{\mathrm{NL,RDCC}}(k,a)
    = \Phi_{\mathrm{NL}}(k,a)
      \left[
        1 - \chi_{\mathrm{NL}}
        \left(\frac{k}{k_{\mathrm{rel}}}\right)^{\alpha}
      \right].
\end{equation}

The nonlinear mode-coupling kernel becomes
\begin{equation}
    F_{2,\mathrm{RDCC}}(\mathbf{k}_{1},\mathbf{k}_{2})
    = F_{2}(\mathbf{k}_{1},\mathbf{k}_{2})
      \left[
        1 - \chi_{F}
        \left(
          \frac{k_{1} + k_{2}}{2k_{\mathrm{rel}}}
        \right)^{\alpha}
      \right].
\end{equation}

This modification reflects the scale-dependent evolution of the gravitational 
potential.
% ---------------------------------------------------------
\section{The Weak-Lensing Bispectrum}

The convergence bispectrum is defined as
\begin{equation}
    B_{\kappa}(\ell_{1},\ell_{2},\ell_{3})
    = \int_{0}^{\chi_{H}}
      d\chi\,
      \frac{
        W^{3}(\chi)
      }{
        \chi^{4}
      }
      B_{m}\left(
        \frac{\ell_{1}}{\chi},
        \frac{\ell_{2}}{\chi},
        \frac{\ell_{3}}{\chi};
        z(\chi)
      \right).
\end{equation}

In RDCC, the matter bispectrum becomes
\begin{equation}
    B_{m,\mathrm{RDCC}}
    = B_{m}
      \left[
        1 - \chi_{B}
        \left(
          \frac{k_{1} + k_{2} + k_{3}}{3k_{\mathrm{rel}}}
        \right)^{\alpha}
      \right].
\end{equation}

The convergence bispectrum becomes
\begin{equation}
    B_{\kappa,\mathrm{RDCC}}
    = B_{\kappa}
      \left[
        1 - \chi_{\kappa}
        \left(
          \frac{\ell_{1} + \ell_{2} + \ell_{3}}{3\ell_{\mathrm{rel}}}
        \right)^{\alpha}
      \right].
\end{equation}
% ---------------------------------------------------------
\section{Configuration Dependence and RDCC Signatures}

The relaxation mechanism enhances the coherence of large-scale modes. The 
weak-lensing bispectrum therefore exhibits:

\begin{itemize}
    \item enhanced equilateral configurations at large angular scales,
    \item suppressed squeezed configurations at small scales,
    \item a characteristic turnover near $\ell_{\mathrm{rel}}$,
    \item a modified redshift dependence of the bispectrum amplitude.
\end{itemize}

These signatures provide a direct observational probe of the RDCC gravitational 
field beyond the two-point level.
% ---------------------------------------------------------
\section{Redshift Evolution of the RDCC Bispectrum}

The redshift evolution of the RDCC bispectrum is given by
\begin{equation}
    \frac{\partial B_{\kappa,\mathrm{RDCC}}}{\partial z}
    = \frac{\partial B_{\kappa}}{\partial z}
      \left[
        1 - \chi_{\kappa}
        \left(
          \frac{\ell_{1} + \ell_{2} + \ell_{3}}{3\ell_{\mathrm{rel}}}
        \right)^{\alpha}
      \right]
      - B_{\kappa}\,
        \chi_{\kappa}\,\alpha\,
        \left(
          \frac{\ell_{1} + \ell_{2} + \ell_{3}}{3\ell_{\mathrm{rel}}}
        \right)^{\alpha}
        \frac{\partial \ln \ell_{\mathrm{rel}}}{\partial z}.
\end{equation}

This term is unique to RDCC and provides a direct probe of the relaxation scale.
% ---------------------------------------------------------
\section{Conclusion}

The weak-lensing bispectrum in the Relaxation-Driven Cyclic Cosmology reflects the 
scale-dependent evolution of the gravitational potential and the nonlinear matter 
distribution. The relaxation mechanism modifies the nonlinear coupling, the 
potential depth, and the three-point structure of the shear field in a 
scale-dependent manner, producing distinctive signatures in cosmic shear surveys. 
These effects provide a direct observational window into the relaxation scale and 
the temporal structure of the RDCC gravitational field. The present Companion 
establishes the theoretical foundation for future studies of non-Gaussianity, 
matter evolution, and the interplay between light and gravity in RDCC.

\bibliographystyle{apsrev4-2}
\nocite{*}
\bibliography{references}

\end{document}
